
\lesson{15/10}
The Dirac Lagrangian admits another continuous internal symmetry, the
\textit{chiral} phase transformation
$$\psi\rightarrow\psi' = e^{i\beta\gamma_5}\rightarrow\delta\psi=i\beta\gamma_5$$
and its Noether current turns out to be
$$J^\mu_5 = \bar\psi\gamma^\mu\gamma_5\psi$$
and is called the \textit{axial current}. The associated Noether charge is
$$Q_5 = \int_{}^{}\dd^3x\bar\psi\gamma^0\gamma_5\psi=
\int_{}^{}\dd^3x\left(\psi^\dagger_L\psi_L - \psi^\dagger_R\psi_R\right)$$
Chirality is important for QFT, because it can only be "detected" by weak
interactions: the Weyl representation is convenient to distinguish the left
and right spinor, while the other, the \textbf{Dirac} representation, is useful
to distinguish particles from antiparticles.
\subsection{Properties of the Dirac Equation}
\textbf{Notation:} Given a 4-vector $a^\mu$, it is sometimes useful to use
$$\slashed{a} := \gamma^\mu a_\mu,$$
even for the 4-gradient:
$$\slashed{\pd} = \gamma^\mu\pd_\mu.$$
So the Dirac equation can be written as
$$(i\slashed{\pd}-m)\psi=0.$$
\subsubsection{Covariance}
Consider a Lorentz transformation of the spinor:
$$\psi(x^\mu)\rightarrow\psi'(x'^\mu) = S(\Lambda)\psi(x^\mu)$$
where
$$S(\Lambda) = e^{\flatfrac{-i}{2}\omega_{\rho\sigma}M^{\rho\sigma}}:$$
how does the Dirac equation transform?
$$(i\gamma^\mu\pd_\mu'-m)\psi'(x')$$
the transformation of the derivative is
$$\pd_\mu' = \underbrace{\Lambda\indices{_\mu^\nu}}_{:=\left(
\Lambda\indices{^\nu_\nu}\right)^{-1}}\pd_\nu$$
and the $\gamma^\mu$ is independent of $x^\mu$: the equation becomes
$$(i\gamma^\mu\Lambda\indices{_\mu^\nu}\pd_\nu-m)S(\Lambda)\psi(x^\mu)=0$$
Applying $S^{-1}(\Lambda)$ yields
$$\left(S^{-1}(\Lambda)\gamma^\mu S(\Lambda)\Lambda\indices{_\mu^\nu}\pd_\nu
-m\right)\psi(x^\mu)=0$$
\textbf{Note:} The $S(\Lambda)$ operator and the $\gamma^\mu$ operate in the same
vector space (that of Dirac spinors), while $\Lambda$ and $\pd$ operate on the
spacetime coordinates: this implies that $S(\Lambda)$ does not commute with
$\gamma^\mu$, but it does with $\Lambda$ and $\pd$.

The equation is covariant if and only if
$$S^{-1}(\Lambda)\gamma^\mu S(\Lambda)= \gamma^\rho\Lambda\indices{^\mu_\rho}.$$
Recalling
$$S(\Lambda) = \exp\left[-\frac{i}{2}\omega_{\rho\sigma}M^{\rho\sigma}\right],\;
M^{0i}=-\frac{i}{2}\begin{pmatrix}
    \sigma^i&\mathbb{O}\\
    \mathbb{O}&-\sigma^i
\end{pmatrix},\;
M^{ij}=\frac{1}{2}\varepsilon^{ijk}\begin{pmatrix}
    \sigma_k&\mathbb{O}\\
    \mathbb{O}&\sigma_k
\end{pmatrix},
$$
one can define
$$\sigma^{\mu\nu}=\frac{i}{2}\left[\gamma^\mu,\gamma^\nu\right]$$
so that the Lorentz generators can be expressed as
$$M^{\mu\nu}=\frac{1}{2}\sigma^{\mu\nu};$$
so the operator is
$$S(\Lambda) = e^{\flatfrac{-i}{4}\omega_{\mu\nu}\sigma^{\mu\nu}}$$
\subsubsection{Klein-Gordon strikes back}
Another way to write the equation is by multiplying it by $-i\gamma^\mu\pd_\mu-m$:
\begin{gather*}
(-i\gamma^\mu\pd_\mu-m)(i\gamma^\nu\pd_\nu -m)\psi=0\\
(\gamma^\mu\gamma^\nu\pd_\mu\pd_\nu+m^2)\psi=0\\
\end{gather*}
the $\gamma^\mu\gamma^\nu$ needs to be symmetrized and yields the anticommutator
$$(\frac{1}{2}\underbrace{\{\gamma^\mu\gamma^\nu\}}_{=2g^{\mu\nu}}\pd_\mu\pd_\nu
+m^2)\psi = 0$$
and so the Klein Gordon equation pops out
$$\left(\Box + m^2\right)\psi=0$$
which means that \textbf{each component of the Dirac spinor is a scalar field},
because its equation of motion is the K-G equation.
\subsubsection{Parity invariance}
The reason Dirac spinors were introduced was to construct a parity-invariant
theory: how do these spinors behave under parity transformations?

Consider a parity transformation
$$x^\mu\to x'^\mu = (x^0,-\mathbf{x})$$
which acts on the field as
$$\psi(x^\mu)\to\psi'(x'^\mu) = S_p\psi(x^\mu).$$
The Dirac equation transforms like
$$(i\gamma^0\pd_0-i\gamma^i\pd_i-m)S_p\psi(x^\mu)=0$$
applying $S^{-1}_p$ to the left
$$(i S^{-1}_p\gamma^0S_p\pd_0-iS_p^{-1}\gamma^i S_p\pd_i-m)\psi(x^\mu)$$
This is invariant if and only if
$$S^{-1}_p\gamma^0 S_p = \gamma^0,\; S^{-1}_p\gamma^i S_p = \gamma^i
\longleftrightarrow S_p = \eta_p\gamma^0$$
where $\eta_p$ is a complex constant called \textit{intrinsic parity}.

Since the quantity $\psi^\dagger\psi$ must be invariant under parity (? showed
using Weyl spinors???), the intrinsic parity $\eta_p$ must be so that
$$|\eta_p|^2=1 \leftrightarrow\eta_p=\pm1$$
